\exnumber{07}
\extitle{UI I - Ex 8 - Défi UI - Montres}
\exfiles{ex07/fichier Figma, nomme Name\_FirstName\_DeliverableName\_Date\_VersionNumber}
\exauthorize{Ceux nommés dans l'exercice}
\makeheaderfiles

\textbf{Instructions :}

Copiez ces images sur Figma

Screen1
Screen2
Screen3
Screen4

\textbf{Livrable :}

\begin{itemize}
\item Un fichier Figma avec une copie de l'image et des éléments Figma
\item Vérifiez :
\begin{itemize}
\item Vous utilisez des symboles et des variants
\item Tous vos calques sont nommés et organisés
\item Vous utilisez des overlays
\item Vous utilisez des masques pour changer la forme de l'image importée
\end{itemize}
\item Déposez le fichier Figma sur Github, nommé ainsi : « Name\_FirstName\_DeliverableName\_Date\_VersionNumber ». Exemple : « Doe\_John\_MidFiPrototype\_05242024\_V1 ».
\item Ajoutez un titre dans le fichier.
\item Soignez l'aspect global du fichier. Il doit être clair, simple et lisible. Vous pouvez regarder canva.com, sans surcharger le dessin de détails.
\end{itemize}

\textbf{Pistes :}

\begin{itemize}
\item Polices : SF Compact pour les éléments iOS, Futura STD Condensed pour les signes Nike
\item Le frame pour chaque écran iWatch est watch 44 mm 184 x 224
\item Article Figma sur les Symbols et Variants \url{https://help.figma.com/hc/en-us/articles/360056440594-Create-and-use-variants}
\item Vous pouvez utiliser des Libraries existantes en parcourant les ressources de la communauté Figma
\end{itemize}

